\documentclass[french,11pt]{book}
\usepackage{teach}
\usepackage[french]{babel}
\usepackage{amsmath,amssymb}
\usepackage{array}
\newcommand{\R}{\mathbb{R}}
\begin{document}
\begin{center}
{\Large \textbf{Devoir surveillé -- Probabilité conditionelle}}\\[2mm]
Terminale technologique
\end{center}

\section*{Devoir surveillé}
\begin{enumerate}
  \item Une urne contient des jetons rouges et bleus. 45\% des jetons sont rouges. Parmi les rouges, 30\% portent un numéro pair ; parmi les bleus, 70\% portent un numéro pair. Construire un arbre pondéré.
  \item Calculer la probabilité de tirer un jeton rouge et pair.
  \item Calculer la probabilité de tirer un jeton pair, puis la probabilité qu'il soit rouge sachant qu'il est pair.
\end{enumerate}

\end{document}
